\subsection{But\ldots{} risks, and the answers to them}
\label{sec:risks-and-regulation}
\slidesrc{Thinking machines: 80--87}

The slides now switch tone (``But\ldots''): AI systems already take or support decisions
that affect people, and things can go wrong.

\subsubsection{Things that went wrong, or might}
\slidesrc{Thinking machines: 80--84}

\begin{center}
\small
\begin{tabularx}{\linewidth}{@{}l >{\raggedright\arraybackslash}X >{\raggedright\arraybackslash}X@{}}
\toprule
\textbf{Case} & \textbf{What happened} & \textbf{The underlying risk} \\
\midrule
State v.\ Loomis (2016--17) & The Wisconsin Supreme Court allowed judges to use an
  algorithmic risk-assessment score (predicting the risk of re-offending) in sentencing,
  but required a written warning about its limits. The defendant could not inspect how the
  proprietary algorithm worked. & opaque decisions about people's freedom; possible bias;
  no way to contest the reasoning \\
``Smart courts'' in China (2019) & Courts adopt AI assistants (e.g.\ Shanghai's
  ``System 206'') to check evidence and support judges. & automation of justice; who is
  accountable? \\
ChatGPT banned in Italy (2023) & The Italian data protection authority temporarily
  blocked ChatGPT over privacy concerns (personal data used for training, no age checks). &
  privacy and data protection \\
AI leaders at the UN Security Council & The heads of OpenAI (Sam Altman) and Anthropic
  (Dario Amodei) gave separate briefings on AI safety to the UN Security Council. &
  global safety and security \\
Incidents with AI agents (slide 84) & News reports about security incidents involving AI
  platforms and AI agents that were misused or hijacked. & autonomous agents acting in the
  real world can be attacked or misbehave \\
\bottomrule
\end{tabularx}
\captionof{table}{The cautionary examples of slides 80--84.}
\end{center}

\begin{nbintuition}[Why symbolic AI helps here]
In State v.\ Loomis the core problem is that nobody could \emph{see} why the system gave a
certain score. A symbolic system (a set of logical rules, a Prolog program, a BDI agent
with explicit plans) can \emph{explain} its conclusions: ``I concluded $X$ because of rule
$R$ applied to facts $F_1$ and $F_2$''. Explainability, verifiability and accountability
are among the strongest arguments for the symbolic side of AI.
\end{nbintuition}

\subsubsection{And hence: principles and guidelines}
\slidesrc{Thinking machines: 85--87}

\begin{itemize}
  \item \textbf{Asilomar AI Principles} (Future of Life Institute, 2017): 23 principles
        developed at the 2017 Asilomar conference, grouped into research issues, ethics and
        values (e.g.\ safety, transparency, responsibility, human control, privacy), and
        longer-term issues.
  \item \textbf{Ethics Guidelines for Trustworthy AI} (European Commission, High-Level
        Expert Group on AI, 8 April 2019). Trustworthy AI must be \emph{lawful},
        \emph{ethical} and \emph{robust}, and should meet seven key requirements (listed
        below, not in the slides).
  \item \textbf{Calls to slow down} (New York Times, September 2026): top AI leaders, among
        them the head of Anthropic, argued that capabilities are advancing so quickly that
        greater safety controls are needed across the industry.
\end{itemize}

\begin{nbfact}{The seven requirements for trustworthy AI (EU, 2019)}
\notinslides
\begin{enumerate}[nosep]
  \item human agency and oversight;
  \item technical robustness and safety;
  \item privacy and data governance;
  \item transparency (including explainability);
  \item diversity, non-discrimination and fairness;
  \item societal and environmental well-being;
  \item accountability.
\end{enumerate}
These guidelines were followed by the EU \emph{AI Act} (in force since 2024), which
regulates AI systems according to their level of risk.
\end{nbfact}
